\chapter{综合实例练习：代码发布、智能体编程与 PyTorch}
\label{cp:ms2-practice}

\section{代码打包与发布（Packaging and Shipping Code）}

\begin{exercise}{用 venv 创建并激活隔离的虚拟环境}
\textbf{要求}：课程指出，不同项目的依赖必须相互隔离，避免版本冲突。使用标准库 \code{venv} 为本项目建立独立虚拟环境，激活后确认解释器位置与版本。

\textbf{操作命令}：
\begin{lstlisting}
python3.14 -m venv .venv
source .venv/bin/activate
which python
python --version
\end{lstlisting}

\textbf{结果说明}：如图\ref{fig:ms2-01}所示，激活后 shell 提示符前出现 \code{(.venv)} 前缀；\code{which python} 指向项目目录 \code{.venv/bin/python} 而非系统解释器，\code{python --version} 输出 Python 3.14.7，说明后续所有 \code{pip} 安装都只作用于该项目环境。

\begin{figure}[H]
    \centering
    \reportimage[width=0.95\textwidth]{shots/01_venv.png}
    \caption{创建并激活 venv 虚拟环境}
    \label{fig:ms2-01}
\end{figure}
\end{exercise}

\begin{exercise}{用 pip 安装依赖并观察传递依赖}
\textbf{要求}：课程讲解了依赖解析与传递依赖（transitive dependencies）。安装第三方库 \code{requests}，观察 pip 自动安装的间接依赖，并用 \code{pip list} 查看环境内全部包。

\textbf{操作命令}：
\begin{lstlisting}
pip install requests
pip list
python -c 'import requests; print(requests.__version__)'
\end{lstlisting}

\textbf{结果说明}：如图\ref{fig:ms2-02}所示，pip 除了安装 requests 2.34.2 本身，还自动解析并安装了它依赖的 certifi、charset-normalizer、idna、urllib3 四个传递依赖；\code{pip list} 列出隔离环境内的全部 6 个包，最后导入 requests 并打印版本号，验证安装成功。

\begin{figure}[H]
    \centering
    \reportimage[width=0.92\textwidth]{shots/02_pip.png}
    \caption{pip 安装 requests 及其传递依赖}
    \label{fig:ms2-02}
\end{figure}
\end{exercise}

\begin{exercise}{用 pyproject.toml 与 build 构建 wheel 包}
\textbf{要求}：现代 Python 包通过 \code{pyproject.toml} 声明项目元数据与构建后端，执行 \code{python -m build} 可同时构建源码分发包 sdist（\code{.tar.gz}）和预编译分发包 wheel（\code{.whl}）。wheel 本质上是一个 zip 压缩包；在 \code{[project.scripts]} 中声明的入口会在安装后成为可执行命令。

\textbf{包源码 \code{greeting.py}}（函数均带类型注解，并用 argparse 接收参数）：
\begin{lstlisting}[language=Python]
def greet(name: str) -> str:
    return f"Hello, {name}!"

def main() -> None:
    import argparse
    parser = argparse.ArgumentParser(description="Greet someone")
    parser.add_argument("name", help="要打招呼的对象")
    args = parser.parse_args()
    print(greet(args.name))
\end{lstlisting}

\textbf{清单文件 \code{pyproject.toml}}：
\begin{lstlisting}
[project]
name = "greeting"
version = "0.1.0"
requires-python = ">=3.10"

[project.scripts]
greet = "greeting:main"

[build-system]
requires = ["setuptools>=61.0"]
build-backend = "setuptools.build_meta"
\end{lstlisting}

\textbf{构建命令}（\code{-q} 为安静模式，只保留关键阶段输出）：
\begin{lstlisting}
python -m build -q
ls -l dist
\end{lstlisting}

如图\ref{fig:ms2-03a}所示，构建依次完成 sdist 与 wheel，最终提示 \code{Successfully built}，\code{dist} 目录下同时生成了 \code{greeting-0.1.0.tar.gz} 和 \code{greeting-0.1.0-py3-none-any.whl} 两个产物。

\begin{figure}[H]
    \centering
    \reportimage[width=0.60\textwidth]{shots/03a_build.png}
    \caption{构建 sdist 与 wheel 产物}
    \label{fig:ms2-03a}
\end{figure}

\textbf{查看 wheel 内容、安装并运行入口命令}：
\begin{lstlisting}
unzip -l dist/greeting-0.1.0-py3-none-any.whl
pip install -q ./dist/greeting-0.1.0-py3-none-any.whl
greet World
greet Alice
\end{lstlisting}

如图\ref{fig:ms2-03b}所示，wheel 内除了 \code{greeting.py}，还有 \code{METADATA}、\code{entry\_points.txt}、\code{RECORD} 等元数据文件，共 6 个文件；本地安装后，\code{[project.scripts]} 声明的 \code{greet} 命令可直接在终端运行，分别输出 \code{Hello, World!} 与 \code{Hello, Alice!}，完成了「源码 $\rightarrow$ 构建 $\rightarrow$ 分发包 $\rightarrow$ 安装运行」的完整发布闭环。

\begin{figure}[H]
    \centering
    \reportimage[width=0.72\textwidth]{shots/03b_wheel.png}
    \caption{查看 wheel 内容、安装并运行 greet 命令}
    \label{fig:ms2-03b}
\end{figure}
\end{exercise}

\begin{exercise}{用 pip freeze 锁定依赖版本以复现环境}
\textbf{要求}：课程强调环境可复现性：把当前环境中每个依赖的精确版本导出为锁定文件，其他机器即可据此还原出完全一致的环境。

\textbf{操作命令}：
\begin{lstlisting}
pip freeze > requirements.txt
wc -l requirements.txt
cat requirements.txt
\end{lstlisting}

\textbf{结果说明}：如图\ref{fig:ms2-04}所示，\code{pip freeze} 共导出 24 个包，每一行都以 \code{包名==精确版本} 的形式固定（如 \code{torch==2.14.0}、\code{requests==2.34.2}），连传递依赖也一并锁定。在新环境中执行 \code{pip install -r requirements.txt} 即可复现同一套依赖。

\begin{figure}[H]
    \centering
    \reportimage[width=0.40\textwidth]{shots/04_freeze.png}
    \caption{pip freeze 生成版本锁定文件}
    \label{fig:ms2-04}
\end{figure}
\end{exercise}

\section{智能体编程（Agentic Coding）}

\begin{exercise}{为编程代理编写 AGENTS.md 说明文件}
\textbf{要求}：课程介绍，编程代理启动时会自动读取项目根目录下的 \code{AGENTS.md}（或 \code{CLAUDE.md}），把项目环境与工作约定预载进上下文，从而无需每次重复交代。用 heredoc 创建该文件并回显确认。

\textbf{操作命令}：
\begin{lstlisting}
cat > AGENTS.md <<'EOF'
# AGENTS.md: 给编程代理的项目说明（启动时自动加载到上下文）
## 环境
# Python 3.14，依赖装在 .venv，先 source .venv/bin/activate
## 工作规则
# 1. 修改任何 .py 文件后必须运行 mypy 类型检查
# 2. 所有函数必须带类型注解
# 3. 打包验证命令：python -m build
EOF
cat AGENTS.md
\end{lstlisting}

\textbf{结果说明}：如图\ref{fig:ms2-05}所示，文件中写明了解释器与虚拟环境位置，以及「改完即做类型检查、函数必须带类型注解、打包验证命令」三条规则；\code{cat} 回显确认内容已正确写入。

\begin{figure}[H]
    \centering
    \reportimage[width=0.66\textwidth]{shots/05_agentsmd.png}
    \caption{创建 AGENTS.md 并回显确认}
    \label{fig:ms2-05}
\end{figure}
\end{exercise}

\begin{exercise}{把一次性脚本改写成带 argparse 的正规命令行程序}
\textbf{要求}：课程演示了把随手脚本逐步改造成「可被人和代理可靠调用」的正规 CLI：用 \code{argparse} 声明位置参数与选项、自动生成帮助信息，并为函数补充类型注解。编写 \code{extract\_links.py} 从 Markdown 文件提取链接，支持 \code{--absolute} 选项只保留绝对 URL。

\textbf{核心代码}：
\begin{lstlisting}[language=Python]
def extract(content: str) -> list[str]:
    return LINK_PATTERN.findall(content)

def main() -> None:
    parser = argparse.ArgumentParser(description="从 Markdown 文件提取链接")
    parser.add_argument("path", type=Path, help="Markdown 文件路径")
    parser.add_argument("--absolute", action="store_true", help="只保留绝对 URL")
    args = parser.parse_args()
    content: str = args.path.read_text(encoding="utf-8")
    links = extract(content)
    if args.absolute:
        links = [url for url in links if url.startswith("http")]
    for link in links:
        print(link)
\end{lstlisting}

\textbf{运行命令}：
\begin{lstlisting}
python extract_links.py --help
python extract_links.py sample.md
python extract_links.py sample.md --absolute
\end{lstlisting}

\textbf{结果说明}：如图\ref{fig:ms2-06}所示，\code{--help} 自动生成了用法、位置参数 \code{path} 和选项 \code{--absolute} 的说明；默认运行提取出 3 条链接（含 1 条相对链接 \code{./local.md}），加上 \code{--absolute} 后只保留 2 条 \code{http} 开头的绝对链接，过滤行为符合预期。

\begin{figure}[H]
    \centering
    \reportimage[width=0.62\textwidth]{shots/06_argparse.png}
    \caption{argparse 命令行程序的帮助与两种运行方式}
    \label{fig:ms2-06}
\end{figure}
\end{exercise}

\begin{exercise}{用 mypy 进行静态类型检查}
\textbf{要求}：课程强调「类型检查反馈环」——每次改完代码就让静态检查器代理自动发现类型问题。对上一实例中带完整类型注解的脚本运行 mypy。

\textbf{操作命令}：
\begin{lstlisting}
mypy extract_links.py
\end{lstlisting}

\textbf{结果说明}：如图\ref{fig:ms2-07}所示，mypy 输出 \code{Success: no issues found in 1 source file}，说明脚本的类型注解自洽、没有类型错误，这正是 AGENTS.md 中「修改后必须运行 mypy」规则的一次实际执行。

\begin{figure}[H]
    \centering
    \reportimage[width=0.50\textwidth]{shots/07_mypy.png}
    \caption{mypy 静态类型检查通过}
    \label{fig:ms2-07}
\end{figure}
\end{exercise}

\section{PyTorch 入门实例}

\begin{exercise}{张量的创建、属性与基本运算}
\textbf{要求}：对应教程「开始使用 PyTorch」，创建张量并查看 \code{shape}、\code{dtype}、\code{device} 属性，练习特殊张量、逐元素运算与矩阵乘法。

\textbf{操作命令}：
\begin{lstlisting}
python tensors.py
\end{lstlisting}

\textbf{结果说明}：如图\ref{fig:ms2-08}所示，从列表创建的 2$\times$2 浮点张量属性为 \code{shape: (2, 2)、dtype: torch.float32、device: cpu}；\code{torch.zeros(2,3)} 生成全零张量，\code{x + 10} 对每个元素加 10（广播），\code{x.sum()} 归约为标量 10.0，\code{x @ I} 与单位矩阵相乘后仍等于 x。

\begin{figure}[H]
    \centering
    \reportimage[width=0.58\textwidth]{shots/08_tensors.png}
    \caption{PyTorch 张量创建、属性与运算}
    \label{fig:ms2-08}
\end{figure}
\end{exercise}

\begin{exercise}{requires\_grad 与 backward 自动微分}
\textbf{要求}：对应 autograd 教程：对需要求导的叶子张量设置 \code{requires\_grad=True}，前向计算后调用 \code{backward()}，梯度自动累积到 \code{.grad}，并与手算结果对照。取 $y=wx+b$、$\mathrm{loss}=y^2$，其中 $w=2,\ x=3,\ b=1$。

\textbf{操作命令}：
\begin{lstlisting}
python autograd_demo.py
\end{lstlisting}

\textbf{结果说明}：如图\ref{fig:ms2-09}所示，前向得到 $y=7$、$\mathrm{loss}=49$；反向传播后 $\partial\mathrm{loss}/\partial x=28$、$\partial\mathrm{loss}/\partial w=42$、$\partial\mathrm{loss}/\partial b=14$。以 $x$ 为例手算 $\partial\mathrm{loss}/\partial x=2y\cdot w=2\times7\times2=28$，与自动微分结果完全一致。

\begin{figure}[H]
    \centering
    \reportimage[width=0.52\textwidth]{shots/09_autograd.png}
    \caption{autograd 自动微分并与手算对照}
    \label{fig:ms2-09}
\end{figure}
\end{exercise}

\begin{exercise}{用 torch.nn 定义并训练一个神经网络}
\textbf{要求}：对应教程「使用 torch.nn 创建并训练神经网络」：用 \code{nn.Sequential} 堆叠线性层与 ReLU，搭配均方误差损失和 SGD 优化器，按「清零梯度 $\rightarrow$ 前向 $\rightarrow$ 反向 $\rightarrow$ 更新参数」的训练循环拟合目标函数 $y=3x-2$。

\textbf{核心代码}：
\begin{lstlisting}[language=Python]
model = nn.Sequential(nn.Linear(1, 16), nn.ReLU(), nn.Linear(16, 1))
loss_fn = nn.MSELoss()
optimizer = torch.optim.SGD(model.parameters(), lr=0.05)

for step in range(501):
    optimizer.zero_grad()
    y_pred = model(X)
    loss = loss_fn(y_pred, y_true)
    loss.backward()
    optimizer.step()
\end{lstlisting}

\textbf{操作命令}：
\begin{lstlisting}
python nn_train.py
\end{lstlisting}

\textbf{结果说明}：如图\ref{fig:ms2-10}所示，每 100 步打印一次损失，loss 从初始的 7.230317 持续下降到 0.000413；训练结束后对 $x=0.5$ 预测得到 $-0.4964$，与目标值 $3\times0.5-2=-0.5$ 几乎一致，说明网络成功学到了线性关系。

\begin{figure}[H]
    \centering
    \reportimage[width=0.46\textwidth]{shots/10_nn_train.png}
    \caption{nn.Sequential 网络的训练损失与预测结果}
    \label{fig:ms2-10}
\end{figure}
\end{exercise}

\section{本章小结}

本章 10 个实例覆盖了三条主线：其一，代码发布侧实践了虚拟环境隔离、pip 依赖解析与传递依赖、基于 \code{pyproject.toml} 的 wheel 构建—安装—运行闭环，以及用 \code{pip freeze} 锁定可复现环境；其二，智能体编程侧实践了用 \code{AGENTS.md} 向代理预置项目约定、用 argparse 与类型注解把脚本工程化、用 mypy 建立自动类型检查反馈环；其三，PyTorch 侧完成了张量运算、autograd 自动微分（与手算互相验证）和 \code{torch.nn} 完整训练循环。所有操作均在终端逐条执行并保留了命令—输出截图，实验过程可按本章命令完整复现。
